\clearpage
\chapter{读谱与记谱}

\classquote{不积跬步，无以至千里；不积小流，无以成江海。}{《荀子·劝学》}{读谱没有捷径：每天认一小节音高与时值，积少成多，谱面才会变成手上的动作。先求准确，再求速度；先能把谱面翻译成卷帘，才谈得上视奏。}

\label{sec:reading-notation}

\section*{学完这一章能做什么}
\begin{itemize}
  \item 看出一张五线谱的谱号、拍号、调号，并说出某条线或某个间对应哪个音。
  \item 把常见音符、休止、附点与连音换算成拍数和时值，不再靠“看起来差不多”猜。
  \item 认出力度、速度与常用表情记号，知道它们在演奏里大概意味着什么。
  \item 把一段简单谱面翻译成音名、MIDI 编号与钢琴卷帘行号，手动抄进夜半。
\end{itemize}

\section{五线谱：线和间不是装饰}

五线谱由五条平行横线组成，线与线之间叫“间”。从下往上数，五条线、四个间各代表一个音级；相邻的线或间相差一个自然音级，而不是固定差两个半音。谱号决定这套位置从哪个音开始，所以同一根线在高音谱号与低音谱号里是完全不同的音。

高音谱号（treble clef）最常用。它的五条线从下到上依次是 E4、G4、B4、D5、F5，四个间从下到上是 F4、A4、C5、E5。一个简单的记法是“下到上 E G B D F”，间是 F A C E。中央 C（C4）不在五线之内，而在高音谱表下方加一条短横线（加线）上；D4 在中央 C 上方的间里，E4 回到第一线。

低音谱号（bass clef）的五条线从下到上是 G2、B2、D3、F3、A3，四个间是 A2、C3、E3、G3。中央 C 在低音谱表上方第一条加线上。钢琴大谱表把高音谱表与低音谱表上下并排，中央 C 在两张谱表之间；看到大谱表时，先看音符靠哪一边，再判断用哪个谱号读。

\begin{center}
\begin{minipage}{0.96\linewidth}
\centering
\begin{tikzpicture}[x=1cm,y=1cm,font=\scriptsize]
  \foreach \y in {0,1,2,3,4} {\draw[gray] (0,\y) -- (10,\y);}
  \draw[gray] (0,-1) -- (10,-1);
  \node[left,font=\tiny] at (0,4) {第五线};
  \node[left,font=\tiny] at (0,0) {第一线};
  \node[left,font=\tiny] at (0,-1) {加线};
  \node[below,font=\tiny] at (0.3,-1.55) {高音谱号};
  \foreach \x/\y/\n in {1.2/-1/C4,2.4/-0.5/D4,3.6/0/E4,4.8/0.5/F4,6.0/1/G4,7.2/1.5/A4,8.4/2/B4,9.3/2.5/C5} {
    \fill (\x,\y) ellipse (0.10 and 0.07);
    \node[above=1pt,font=\tiny] at (\x,\y) {\n};
  }
\end{tikzpicture}

图 9-1：高音谱表的线与间，以及中央 C 的加线位置。相邻线/间是一个自然音级；E--F、B--C 只差一个半音，其他相邻音级差两个半音。
\end{minipage}
\end{center}

读谱时不要从第一个音开始逐个“数格子”。先找锚点：高音谱号的中央 C、第一线 E4、第三线 B4；低音谱号的中央 C、第一线 G2、第三线 D3。锚点定住后，其他音按相邻线/间推算，速度会快很多。

\section{音符、休止与时值}

音符的“形状”表示时值，不表示音高。全音符是一个空心椭圆，二分音符加符干，四分音符实心加符干，八分音符再加一条符尾，十六分音符两条符尾。符尾越多，时值越短；每多一条符尾，时值减半。

在 4/4 拍里，一个全音符等于四拍，二分音符两拍，四分音符一拍，八分音符半拍，十六分音符四分之一拍。附点把一个音符的时值延长它的一半：附点二分音符等于三拍，附点四分音符等于一拍半。连音线（tie）把相邻的同音音符连成一个更长的音，不重新起音；延音线跨小节时最常用。

休止符表示“这里不发声”，时值与同名音符对应：全休止、二分休止、四分休止、八分休止、十六分休止。休止不是“漏写”，它是节奏的一部分；一段音乐如果没有休止，听感会一直顶着，缺少呼吸。

\begin{center}
\begin{minipage}{0.94\linewidth}
\centering
\begin{tikzpicture}[x=1cm,y=1cm,font=\scriptsize]
  \foreach \r/\lab/\w in {0/全音符/4,1/二分音符/2,2/四分音符/1,3/八分音符/0.5} {
    \draw[fill=gray!10] (0,-\r*0.8) rectangle ++(\w,0.45);
    \node[left,font=\tiny] at (0,-\r*0.8+0.22) {\lab};
    \node[right,font=\tiny] at (\w,-\r*0.8+0.22) {\w 拍};
  }
  \foreach \x in {0,1,2,3,4} {\draw[gray!35] (\x,0.1) -- (\x,-2.5);}
  \node[below,font=\tiny] at (2,-2.55) {4/4 拍，一小节四拍};
\end{tikzpicture}

图 9-2：四种常见时值在 4/4 拍里的长度比例。长度按拍数等比，不按音符图形大小。
\end{minipage}
\end{center}

\begin{center}
\small
\begin{tabular}{p{0.16\linewidth}p{0.18\linewidth}p{0.16\linewidth}p{0.32\linewidth}}
\toprule
音符 & 4/4 拍数 & 休止符 & 说明 \\
\midrule
全音符 & 4 & 全休止 & 一小节的基本单位 \\
二分音符 & 2 & 二分休止 & 两个二分音符填满一小节 \\
四分音符 & 1 & 四分休止 & 最常见的拍单位 \\
八分音符 & 0.5 & 八分休止 & 两个八分音符等于一个四分 \\
十六分音符 & 0.25 & 十六分休止 & 四个十六分等于一个四分 \\
附点四分 & 1.5 & --- & 延长一半，常与八分组合 \\
\bottomrule
\end{tabular}

表 9-1：常见音符、休止与 4/4 拍时值。附点与连音会改变实际长度，读谱时要单独算。
\end{center}

拍号写在谱表开头：上面的数字是每小节拍数，下面的数字是拍的单位。4/4 表示每小节四拍、四分音符为一拍；6/8 表示每小节六拍、八分音符为一拍。拍号不是速度；速度由 BPM 或速度记号决定。4/4 的 120 BPM 表示每分钟 120 个四分音符，6/8 的 120 BPM 则表示每分钟 120 个八分音符。

\section{调号、力度与表情}

调号写在拍号旁边，用升号或降号告诉你“哪些音要升高/降低”。一个升号通常是 F\sharpsym，对应 G 大调或 E 小调；一个降号通常是 B\flatsym，对应 F 大调或 D 小调。调号只改那些被标记的音级，不会把整条音阶的所有音都升高；具体音阶与调式的关系见第 4 章。

力度记号写在音符下方或谱表之间。$p$（piano）是弱，$mp$（mezzo-piano）是中弱，$mf$（mezzo-forte）是中强，$f$（forte）是强；$<$ 形的渐强（crescendo）与 $>$ 形的渐弱（diminuendo）表示在一段内连续变化。力度不是“音量旋钮的绝对值”，而是相对同一段音乐的强弱层次；同样标 $mf$，在安静段与激烈段听到的实际响度可以不同。

表情与发音记号决定“怎么发出这个音”。连音线（slur）表示一段连贯演奏，断奏（staccato）表示短促、断开，重音（accent）表示这个音要更突出，延长记号（fermata）表示在一个音或休止上停留。它们没有统一的毫秒数，演奏者要根据风格处理；读谱时先认出符号，再把意图转成你的时值、力度与音色。

\begin{center}
\small
\begin{tabular}{p{0.18\linewidth}p{0.22\linewidth}p{0.46\linewidth}}
\toprule
记号 & 名称 & 读谱时先做什么 \\
\midrule
$p$ / $mf$ / $f$ & 弱 / 中强 / 强 & 定这一段的动态层级，不要逐音跳变 \\
$<$ / $>$ & 渐强 / 渐弱 & 在指定段内连续变化，不要只在最后一个音爆发 \\
$\bullet$ 符头加小点 & 断奏 staccato & 把音符时值缩短，留出空隙 \\
符头上方短横 & 保持/连奏 tenuto & 把音奏满，甚至略长 \\
符头上方尖角 & 重音 accent & 加强这一个音的起音与力度 \\
弧线跨多个音 & 连音线 slur & 连贯演奏，音与音之间不重新起音 \\
符头上方半圆点 & 延长 fermata & 在该音或休止上停留，长度由指挥/演奏者定 \\
\bottomrule
\end{tabular}

表 9-2：常见力度与表情记号。先认符号，再决定时值、力度与音色，不要只记名字。
\end{center}

\section{简谱与 TAB 的存在意义}

简谱用 1--7 表示七个自然音级，0 表示休止；数字上方或下方的点表示高八度或低八度，数字下方的横线表示时值减半。它适合歌唱教学、民乐与快速记谱：读起来直接对应唱名，转调时也容易整体移动。它的弱点是缺少五线谱的固定音高位置与和声排列信息，遇到转位、密集和声或多声部时，读起来不如五线谱直观。

TAB（吉他谱）用弦与品位表示音，六条线对应六根弦，线上的数字是品位。它适合吉他/贝斯演奏，能直接告诉演奏者“按哪里”，但不直接告诉你音名、音区与和声功能；不同定弦的乐器也不能共用同一份 TAB。TAB 不是“更简单的五线谱”，它是另一种演奏坐标。

三种记谱法没有谁替代谁：五线谱擅长固定音高与多声部，简谱擅长唱名与教学，TAB 擅长具体乐器指法。写歌时可以用简谱记旋律草稿，用五线谱检查和声，用 TAB 给吉他手排练；到了夜半，最终都要落到音高行与时值上。

\section{把谱面翻译成钢琴卷帘}

夜半当前没有内建五线谱视图，也没有可用的 MusicXML 导入入口；读谱仍要在纸面或外部谱面上完成，再把音高、时值、力度手动抄进钢琴卷帘。翻译顺序如下。

\begin{enumerate}
  \item 先看谱号，确定五条线各是什么音；再看调号，把该升/该降的音记住；最后看拍号，确定每小节几拍。
  \item 把音符位置翻成音名，再翻成 MIDI 编号：C4 = 60，D4 = 62，E4 = 64，F4 = 65，G4 = 67，A4 = 69，B4 = 71，C5 = 72。
  \item 注意五线谱的相邻线/间是自然音级，不是固定半音。E--F 与 B--C 只差一个半音，在卷帘里只隔一行；C--D、D--E、F--G、G--A、A--B 差两个半音，在卷帘里隔两行。抄谱时按音名查 MIDI 编号，不要按“线数×2”硬推。
  \item 把时值换成拍数，再换成卷帘的格子。4/4 拍、120 BPM 时，一个四分音符是一拍；卷帘的网格若按十六分音符细分，一个四分音符占四格。附点、连音与跨小节延音要单独加长，不能只看音符图形。
  \item 力度记号转成 velocity（力度值），渐强/渐弱写成一段内逐步变化的力度；断奏把音符时值缩短并留空隙，连音线让相邻音重叠或至少不断开。卷帘不会自动识别谱面记号，这些都要手动做。
\end{enumerate}

\screenshot{piano-roll-staff-mapping}{谱面音名与卷帘行号对照}

\textbf{范围说明：}夜半没有五线谱/简谱/TAB 的谱面显示、自动识谱或 MusicXML 导入入口；钢琴卷帘只显示音高行、时值与力度，不显示谱号、调号、力度记号或连音线。读谱的“翻译”步骤必须由你完成；应用不会替你判断某个音是否按调号升降，也不会把 staccato 自动转成短音符。

\section{名曲里的记谱与卷帘对照}

\begin{itemize}
  \item \textbf{Beethoven《欢乐颂》}：全曲主旋律高度集中在连续五个音内，把简谱数字、五线谱线间位置与夜半钢琴卷帘的台阶高度放在一起看，三分钟就能打通三种读法。
  \item \textbf{《小星星》（Twinkle Twinkle Little Star）}：开头的同音重复加纯五度上跳，在五线谱上是直观的“跳两格”，在钢琴卷帘上则是清晰的 7 个半音高度差。
  \item \textbf{邓丽君《月亮代表我的心》}：大量附点与跨拍延音线，在乐谱里靠附点与连线标记，在夜半钢琴卷帘里则直接表现为拉长的音符条长度。
\end{itemize}

\section{练习}


\begin{exercise}
\begin{enumerate}
  \item 高音谱表第一线、第三线、中央 C 分别是什么音？它们的 MIDI 编号是多少？
  \par\noindent 答案：第一线 E4（64），第三线 B4（71），中央 C 是 C4（60，谱表下方加线）。
  \item 4/4 拍里，一个附点二分音符、一个附点四分音符各等于几拍？一小节能放几个附点四分音符？
  \par\noindent 答案：附点二分 3 拍；附点四分 1.5 拍；4/4 一小节最多放两个附点四分（3 拍）再加一个四分（1 拍），不能放三个（4.5 拍超一小节）。
  \item 一个升号的调号通常对应哪两个调？哪个音被升高？
  \par\noindent 答案：G 大调或 E 小调；F 升高为 F\sharpsym。
  \item 把 C4--D4--E4--F4--G4 翻成 MIDI 编号，并说明它们在卷帘里各隔几行。
  \par\noindent 答案：60、62、64、65、67；C--D 隔两行，D--E 隔两行，E--F 隔一行，F--G 隔两行（按半音行计）。
  \item 在夜半里抄一段四音谱例：C4、E4、G4、C5，全部写成四分音符，并给最后一个音更弱的力度；写出你的操作顺序。
  \par\noindent 答案：先建 4/4 工程，打开卷帘；按 C4=60、E4=64、G4=67、C5=72 找行；每音写满一拍的四分音符；选中最后一个音把 velocity 调低；播放确认音高与拍数。
\end{enumerate}
\end{exercise}

\section{帮助主题映射}

\begin{center}
\small
\begin{tabular}{p{0.28\linewidth}p{0.46\linewidth}p{0.18\linewidth}}
\toprule
本章小节 & 应用内帮助锚点 & 备注 \\
\midrule
五线谱与谱号 & \texttt{music-theory-staff} & 图 9-1 \\
音符与休止 & \texttt{music-theory-staff-note-values} & 表 9-1 \\
调号、力度与表情 & \texttt{music-theory-dynamics} & 表 9-2 \\
简谱与 TAB & \texttt{music-theory-jianpu-tab} & 概念 \\
谱面转卷帘 & \texttt{music-theory-piano-roll-notation} & 操作 \\
练习 & \texttt{music-theory-ch9-practice} & 答案在章末 \\
\bottomrule
\end{tabular}
\end{center}

\section*{延伸阅读}
\begin{itemize}
  \item \textbf{Open Music Theory, “Basic notation.”} \url{https://openmusictheory.github.io/basicNotation.html} —— 系统过一遍谱表、谱号、音符与休止，补本章的记谱符号。
  \item \textbf{YouTube, “How to Read Sheet Music from Scratch.”} \url{https://www.youtube.com/watch?v=6aj-mI_J3ww} —— 从谱表与谱号开始示范读谱，适合零基础跟练。
\end{itemize}

